\documentclass[a4paper,12pt]{article}
\usepackage[brazil]{babel}
\usepackage[utf8]{inputenc}
\usepackage{amsmath, amssymb}
\usepackage{geometry}
\geometry{margin=2.5cm}

\begin{document}

\section*{Questão 169 --- ENEM 2024 --- Matemática}

\paragraph{Enunciado:}
A criptografia refere-se à construção e análise de protocolos que impedem terceiros de lerem mensagens privadas.
Júlio César, imperador romano, utilizava um código para proteger as mensagens enviadas a seus generais. Assim,
se a mensagem caísse em mãos inimigas, a informação não poderia ser compreendida. Nesse código, cada letra do
alfabeto era substituída pela letra três posições à frente, ou seja, o ``A'' era substituído pelo ``D'', o ``B'' pelo ``E'', o ``C''
pelo ``F'', e assim sucessivamente.

Qualquer código que tenha um padrão de substituição de letras como o descrito é considerado uma Cifra de César
ou um Código de César. Note que, para decifrar uma Cifra de César, basta descobrir por qual letra o ``A'' foi substituído,
pos isso define todas as demais substituições a serem feitas.

Uma mensagem, em um alfabeto de 26 letras, foi codificada usando uma Cifra de César. Considere a probabilidade
de se descobrir, aleatoriamente, o padrão utilizado nessa codificação, e que uma tentativa frustrada deverá ser eliminada
nas tentativas seguintes.

A probabilidade de se descobrir o padrão dessa Cifra de César apenas na terceira tentativa é dada por:

\bigskip

% Imagem do enunciado (descrição textual)
\noindent\textbf{Imagem:}\
Uma tabela mostrando duas linhas de letras:\
\textbullet\ Linha inferior: Alfabeto original, Letras de A a Z.\
\textbullet\ Linha superior, deslocada três letras para a esquerda a partir do A, mostrando o texto codificado.

\bigskip

\section*{Solução Técnica}
\begin{itemize}
  \item O alfabeto tem 26 letras, porém o deslocamento 0 (sem mudança) não é considerado, então existem 25 padrões possíveis.
  \item Probabilidade de errar na primeira tentativa:
  \[
  P(\text{erro no 1º}) = \frac{24}{25}
  \]
  \item Após eliminar a primeira tentativa, restam 24 opções para a segunda tentativa.
  \item Probabilidade de errar na segunda tentativa:
  \[
  P(\text{erro no 2º} \mid \text{erro no 1º}) = \frac{23}{24}
  \]
  \item Após eliminar a segunda tentativa, restam 23 opções para a terceira tentativa.
  \item Probabilidade de acertar na terceira tentativa:
  \[
  P(\text{acerto no 3º} \mid \text{erros anteriores}) = \frac{1}{23}
  \]
  \item Multiplicando as probabilidades:
  \[
  P(\text{acerto no 3º}) = \frac{24}{25} \times \frac{23}{24} \times \frac{1}{23} = \frac{1}{25}.
  \]
\end{itemize}

\section*{Explicação Didática}
A probabilidade de descobrir o padrão na terceira tentativa significa que as duas primeiras tentativas são erros e somente na terceira ocorre o sucesso.

\begin{enumerate}
  \item Existem 25 padrões possíveis para a Cifra de César (deslocamentos de 1 a 25).
  \item Na primeira tentativa, a probabilidade de errar é $24/25$.
  \item Eliminando essa tentativa errada, restam 24 padrões para a segunda tentativa.
  \item Na segunda tentativa, a probabilidade de errar é $23/24$.
  \item Eliminando essa segunda tentativa errada, restam 23 padrões para a terceira tentativa.
  \item Na terceira tentativa, a probabilidade de acertar é $1/23$.
  \item Portanto, multiplicamos essas probabilidades para obter o evento desejado:
  \[
  \left(\frac{24}{25}\right) \times \left(\frac{23}{24}\right) \times \left(\frac{1}{23}\right) = \frac{1}{25}.
  \]
\end{enumerate}

\section*{Erros Comuns}
\begin{itemize}
  \item Somar probabilidades em vez de multiplicar para eventos sequenciais dependentes.
  \item Não atualizar o espaço amostral ao eliminar uma tentativa errada.
  \item Ignorar a dependência entre as tentativas.
  \item Usar denominadores fixos ao invés de reduzir o espaço amostral progressivamente.
\end{itemize}

\section*{Dicas de Ensino}
\begin{itemize}
  \item Explique como o espaço amostral diminui após cada tentativa eliminada.
  \item Diferencie claramente eventos independentes e dependentes.
  \item Use exercícios para praticar probabilidade condicional.
  \item Mostre diagramas que representem a progressiva redução do número de opções.
  \item Enfatize a aplicação histórica do Código e Cifra de César.
\end{itemize}

\section*{Análise dos Distratores}
\begin{description}
  \item[A] Soma as probabilidades $\tfrac{1}{25} + \tfrac{1}{25} + \tfrac{1}{25}$ ignorando o descarte das tentativas erradas.
  \item[B] Soma frações com denominadores decrescentes $\tfrac{24}{25} + \tfrac{23}{24} + \tfrac{1}{23}$ ao invés de multiplicar.
  \item[C] Multiplica, mas não atualiza corretamente o denominador da segunda tentativa para 24.
  \item[D] Multiplica, porém com denominadores incorretos, misturando valores.
\end{description}

\section*{Perfil da Questão}
\begin{itemize}
  \item \textbf{Habilidade principal}: Probabilidade condicional e eventos dependentes.
  \item \textbf{Habilidades secundárias}: Combinatória, interpretação de texto e noções básicas de criptografia.
  \item \textbf{Nível de dificuldade}: Médio.
  \item \textbf{Requisitos}: Interpretação visual, modelagem e cálculo probabilístico.
  \item \textbf{Observações}: Distratores exploram erros típicos em cálculo de probabilidade condicional e atualização do espaço amostral.
\end{itemize}

\end{document}